\chapter*{Prefacio}
\addcontentsline{toc}{chapter}{Prefacio}
\markboth{Prefacio}{Prefacio}

\section*{Interacción compleja entre dos máquinas simples}
\phantomsection
\addcontentsline{toc}{section}{Interacción compleja entre dos máquinas simples}

Este documento recopila el material técnico extenso —para futuras publicaciones— de los resultados obtenidos y de los desarrollos matemático-físicos realizados entre los años 2023 y 2026 mediante el uso de herramientas de Large Language Models (LLM). De este proceso surgen lo que denominamos, respectivamente, Holografía Modular Triádica y Mecánica Dimensional (HMT–MD). Sin embargo, el origen de este proyecto se remonta a principios del siglo XXI, concretamente al estudio del descubrimiento que subyace en el artefacto inventado por Oumar Haidara Fall en el año 2001. Cabe apuntar aquí que dicho artefacto fue puesto en conocimiento de las diversas autoridades académicas y civiles competentes en los respectivos lugares de origen de ambos autores.

El material acumulado y los conocimientos adquiridos progresivamente a lo largo de más de dos décadas de observación, estudio y análisis, tanto del artefacto como de su fenomenología, constituyeron la base operatoria e informacional que posteriormente aplicamos a los modelos de inteligencia artificial.

Por tanto, el resultado HMT–MD que aquí presentamos, junto con sus derivaciones y demostraciones, viene precedido por una genealogía empírica que, \textit{a priori}, pudiera resultar, en el mejor de los casos, contraintuitiva para el técnico especialista —pues contraviene los postulados y predicciones del marco teórico-interpretativo en el que, por tradición histórica, se inscribe el artefacto, a saber, la mecánica clásica—. Sin embargo, su comprobación metrológica resulta decisiva. Además, salvando prejuicios cognitivos o sesgos de confirmación, se hace directamente demostrable por la vía de los hechos.

Por último, cabría considerar algo que no deja de sorprender como paradoja histórica: cuando empezábamos a razonar con “máquinas pensantes” en “la era de la inteligencia artificial”, tuvimos además la serendipia de desvelar una interacción inédita en la mecánica ordinaria y nos vimos, por ello, obligados a revisar la relación que existe entre dos máquinas simples: la palanca y la rueda.
